% Formatted using the Springer LNCS samplepaper template style.
\documentclass[runningheads]{llncs}

\usepackage[T1]{fontenc}
\usepackage[utf8]{inputenc}
\usepackage{graphicx}
\usepackage{amsmath,amssymb}
\usepackage{url}
\usepackage{tabularx}
\usepackage{booktabs}
\usepackage{float}
\usepackage{array}

\begin{document}

\title{Belief Networks for Fault Diagnosis under Incomplete Information: Applicability and Limitations}
\titlerunning{Belief Networks for Fault Diagnosis}

\author{Pablo Garcia-Legaz Levi\inst{1}}
\authorrunning{P. Garcia-Legaz Levi}

\institute{Vrije Universiteit Amsterdam, Amsterdam, The Netherlands\\
Student number: 2727290}

\maketitle

\begin{abstract}
Fault diagnosis often requires reasoning from observable symptoms to hidden component failures under incomplete information. Bayesian networks are useful for structured diagnostic reasoning, but they require a complete specification of priors and conditional probability tables. This thesis investigates belief networks, represented using Dempster--Shafer belief functions, as a practical alternative when the diagnostic structure is known but some probabilistic information is missing or weakly justified. A modular electrical system is used as a case study. A complete Bayesian network provides a reference structure and validation point, while the Belief Function Machine in MATLAB is used to implement belief-network variants with progressively incomplete information. The results show that when evidence is decisive, the belief network preserves the correct diagnosis, while under missing priors or measurements it exposes unresolved alternatives through belief and plausibility rather than forcing precise posterior rankings. The thesis concludes with practical guidance on when belief networks are recommended for industrial fault diagnosis and where scalability, tooling, and interpretation remain limiting factors.

\keywords{Belief networks \and Dempster--Shafer theory \and Fault diagnosis \and Incomplete information \and Bayesian networks}
\end{abstract}


\section{Introduction}
\label{sec:introduction}

Fault diagnosis is the task of identifying the cause of abnormal behaviour in a technical system \cite{Isermann2006}. Modern industrial machinery, energy systems, and cyber-physical systems are highly interconnected, so a single component failure can propagate through the system and produce symptoms that are not obviously tied to their root cause, with consequences ranging from downtime and financial loss to safety risk \cite{Venkatasubramanian2003PartI,GaoCecatiDing2015PartII}. A diagnostic method must therefore reason from observable symptoms -- alarms, indicator lights, sensor readings, or failed outputs -- back to hidden faults in components, connections, or subsystems.

A minimal example shows the difficulty. Consider a battery, a switch, a cable, and a lamp connected in series, and suppose the lamp does not light. The observation alone does not identify the cause: the battery may be flat, the switch open, the cable broken, or the lamp blown. Additional observations such as indicator LEDs or voltage readings narrow the possibilities but rarely eliminate them. This thesis studies this kind of diagnostic ambiguity using the modular electrical system described in Section~\ref{sec:case_study_diagnostic_model}.

Bayesian networks (BNs) are among the most widely used frameworks for reasoning under uncertainty in structured domains \cite{Pearl1988,Heckerman1995}, and have been applied to fault detection and diagnosis where hidden states must be inferred from imperfect evidence \cite{LernerParrKollerBiswas2000}. A BN encodes faults, internal states, and symptoms as nodes in a directed acyclic graph; before inference can be run, every root node must be assigned a prior probability and every other node a conditional probability table (CPT). Evidence is then propagated to obtain posterior fault probabilities. The limitation relevant here is not that BNs reason poorly, but that they require uncertainty to be expressed as complete and precise probability distributions.

This becomes problematic when the diagnostic structure is known but some probabilities are unavailable, weakly justified, or hard to estimate reliably, for example for rare faults, redesigned components, or parts from a new supplier. Common remedies are to fill the gaps with default assumptions, such as uniform or maximum-entropy values, or to elicit the numbers from domain experts, which can introduce bias and overconfidence \cite{OHagan2006}. Either way, the model still returns precise posteriors that can hide how little the underlying knowledge actually supports them \cite{Shenoy2023IncompleteBN}.

Several uncertainty theories are designed to represent incomplete information directly. Dempster--Shafer, or belief-function, theory assigns support not only to individual hypotheses but also to sets of hypotheses \cite{Dempster1967,Shafer1976}. If the evidence indicates an upstream electrical interruption but does not justify separating a battery fault from an early cable fault, a belief function can commit support to the set \{battery, cable\} as a whole rather than forcing an artificial split. Two quantities are then read from the result: belief, the support committed to a fault, and plausibility, the support compatible with it. Belief and plausibility are treated here as measures of committed and compatible support, not as endpoints of an interval guaranteed to contain a ``true'' probability.

This work builds on Shenoy's method for inference in incomplete Bayesian networks using belief functions \cite{Shenoy2023IncompleteBN}. In that approach, known probabilities are embedded as belief functions while missing ones are left vacuous. The complete BN is therefore not a competitor to be beaten; it is a reference structure and validation point. When all information is present, the belief-function model should reproduce the BN. Once information is deliberately removed, its behaviour under incompleteness can be studied. The aim is not to show that belief networks are more accurate than BNs, but to characterise when they represent genuinely incomplete diagnostic knowledge more transparently.

The thesis is deliberately practical. Rather than proposing new theory, it shows how a practitioner can move from a known diagnostic structure to a belief network, how missing information can be preserved instead of invented, and how the output should be read. Because the motivation is real-world applicability under limited data, the analysis also asks which limitations the approach shares with BNs, which it removes, which it adds, and how close it is to deployment, framed through technology-readiness levels \cite{Mankins1995}. The main research question is:

\begin{quote}
When are belief networks a suitable approach for fault diagnosis under incomplete probabilistic information, and how can such a model be implemented and evaluated in practice?
\end{quote}

It is addressed through five subquestions: how a fault-diagnosis problem can be represented as a belief network (Sections~\ref{sec:case_study_diagnostic_model}--\ref{sec:building_belief_network}); how missing diagnostic information can be represented using belief functions (Section~\ref{sec:representing_missing_information}); how the Belief Function Machine can be used for implementation and validation (Sections~\ref{sec:implementation_bfm} and~\ref{sec:l0_validation}); how belief, plausibility, and rankings behave as information is removed (Section~\ref{sec:experimental_results}); and under what practical conditions the approach is useful or limited (Section~\ref{sec:practical_guidance}).

\section{Background}
\label{sec:background}


\subsection{Fault Diagnosis under Incomplete Information}
\label{sec:fault_diagnosis_uncertainty}

Fault diagnosis links observations to explanations and actions. In engineering systems, the objective is not only to detect abnormal behaviour, but to identify plausible fault causes early enough to support repair, maintenance planning, or risk mitigation \cite{Isermann2006,Venkatasubramanian2003PartI}. Diagnostic methods are commonly grouped into model-based, signal-based, data-driven, and hybrid approaches; in practice, industrial systems often combine physical structure, measured signals, historical data, and expert knowledge \cite{Venkatasubramanian2003PartI,Venkatasubramanian2003PartII,Venkatasubramanian2003PartIII,GaoCecatiDing2015PartII}.

This thesis focuses on model-based graphical diagnosis under incomplete information. The qualitative structure of the system may be known: for example, voltage should propagate from a power supply through control modules to a load, and certain faults interrupt this propagation. However, this structural knowledge does not automatically provide all probabilities needed for probabilistic diagnosis. Observations may be incomplete, several faults may produce similar symptoms, and experts may know that a fault--symptom relation exists without being able to quantify it for every possible condition. The central problem is therefore not only to infer a likely fault, but to represent how incomplete the diagnostic knowledge is.


\subsection{Bayesian Networks as a Reference Framework}
\label{sec:bayesian_networks_reference}

A Bayesian network is a directed acyclic graph together with local probability distributions. Root nodes are assigned priors, while non-root nodes are assigned conditional probability tables. If the network variables are $X_1,\ldots,X_n$ and $Pa_i$ denotes the parents of $X_i$, the joint distribution factorises as
\begin{equation}
P(x_1,\ldots,x_n)=\prod_{i=1}^{n}P(x_i\mid pa_i).
\end{equation}
This factorisation is attractive for diagnosis because faults, internal states, and observable symptoms can be encoded modularly \cite{Pearl1988,Heckerman1995}. Once evidence is observed, Bayesian inference updates the posterior probabilities of the fault variables; BNs have therefore been used for diagnostic problems involving hidden system states and imperfect evidence \cite{LernerParrKollerBiswas2000,VerbertBabuskaDeSchutter2017}.

The relevant limitation is the requirement for precise probabilistic parameters. A complete BN requires all priors and CPT entries to be specified. For rare faults, new components, or unobserved parent-state combinations, those numbers may be unavailable or weakly justified. Completing the model with arbitrary values makes inference possible, but the resulting posterior can appear more precise than the underlying knowledge warrants \cite{Shenoy2023IncompleteBN}. This thesis therefore uses a complete BN as a reference under full information, while evaluating belief networks in cases where such precise quantification is not justified.


\subsection{Dempster--Shafer Theory and Belief Functions}
\label{sec:ds_theory}

Dempster--Shafer theory, also called belief-function theory, represents uncertain evidence using a basic probability assignment (BPA), or mass function \cite{Dempster1967,Shafer1976}. Let $\Theta$ be a finite frame of discernment, i.e. the set of mutually exclusive and exhaustive hypotheses. A BPA is a function
\begin{equation}
m:2^\Theta\rightarrow[0,1]
\end{equation}
such that
\begin{equation}
m(\emptyset)=0, \qquad \sum_{A\subseteq\Theta}m(A)=1.
\end{equation}
The subsets $A\subseteq\Theta$ for which $m(A)>0$ are called focal elements. A focal element is a hypothesis, or set of hypotheses, to which the model assigns direct support. Ordinary probability is a special case: if all focal elements are singletons, such as $\{\theta_i\}$, the mass function behaves like a probability distribution. Belief functions generalise this by allowing mass on larger subsets, which is how they express partial ignorance. Complete ignorance over a variable is represented by the vacuous BPA $m(\Theta)=1$.

From a mass function, belief and plausibility are defined as
\begin{equation}
Bel(A)=\sum_{B\subseteq A}m(B),
\end{equation}
\begin{equation}
Pl(A)=\sum_{B\cap A\neq\emptyset}m(B)=1-Bel(\bar{A}).
\end{equation}
Belief is support committed to $A$; plausibility is support compatible with $A$ \cite{Shafer1976,SmetsKennes1994}. In this thesis, they are used as measures of committed and compatible support rather than as ordinary posterior probabilities. A low belief but high plausibility for a fault means that the fault is not directly supported individually, but remains possible under the available evidence.

Distinct bodies of evidence can be combined using Dempster's rule:
\begin{equation}
(m_1\oplus m_2)(A)=\frac{\sum_{B\cap C=A}m_1(B)m_2(C)}{1-\sum_{B\cap C=\emptyset}m_1(B)m_2(C)}, \quad A\neq\emptyset.
\end{equation}
The rule strengthens support where sources agree and normalises away conflicting mass \cite{Shafer1976}. It must be used carefully because it assumes distinct bodies of evidence; otherwise, the same information may be double-counted. This issue is important in diagnostic systems, where several observations may originate from the same physical dependency or sensor chain. Alternative cautious combination rules have been proposed for non-distinct or overlapping evidence \cite{Denoeux2008,Shenoy2023DistinctBeliefFunctions}.


\subsection{Belief Networks and the BFM Approach}
\label{sec:belief_networks_bfm}

Belief functions can be organised in graphical models. In this thesis, a belief network is a graphical diagnostic model whose local uncertainty relations are represented using belief functions rather than only precise conditional probabilities. Foundational work on local computation showed that belief functions can be propagated modularly in graphical structures \cite{ShaferShenoyMellouli1987}. More recent work defines conditional belief functions for directed graphical models, providing the analogue of CPTs for belief networks \cite{JirousekKratochvilShenoy2023}.

The key source for this thesis is Shenoy's approach to inference in incomplete Bayesian networks \cite{Shenoy2023IncompleteBN}. If a BN is complete, the corresponding belief-function model should reproduce its probabilistic behaviour. If some priors or CPT rows are missing, the known parts can be embedded as belief functions while the missing parts are represented vacuously. The implementation uses the Belief Function Machine (BFM), a MATLAB environment for reasoning with belief functions \cite{GiangShenoy2003BFM}. BFM represents variables, relations, valuations, and conditional valuations in a text-based input language; conditional valuations are converted using Smets' conditional embedding so that they can be combined and marginalised in the belief-function framework.

\section{Case Study and Diagnostic Model}
\label{sec:case_study_diagnostic_model}

This section introduces the diagnostic case study used to demonstrate and evaluate the method. The purpose is not to provide a full engineering specification of the electrical system or a complete listing of all probability tables. Those details belong in the appendix. The main text gives only the information needed to understand the belief-network construction and the experiments.

\subsection{Electrical System and Diagnostic Task}
\label{sec:electrical_system_diagnostic_task}

The case study is based on a modular electrical system whose overall function is to operate a load, represented by a lamp. The system consists of ten modules connected in series: one power supply module, eight control modules, and one load module. The power supply module contains a 12V lithium battery and an indicator LED. Each control module contains a switch and a switch indicator LED. The load module contains the main lamp and a lamp indicator.

The diagnostic task is to infer which hidden fault, or group of faults, explains the observed behaviour of the system. The model distinguishes between hidden fault variables, intermediate voltage variables, and observable variables. Fault variables represent possible root causes, including a power supply short, battery exhaustion, switch detachment faults, inter-module cable faults, and lamp failure. Voltage variables represent propagation of electrical power through the system, with $V_0$ denoting voltage at the power supply output and $V_1$ to $V_8$ denoting post-switch voltage after each control module. Observable variables include the power supply LED, switch indicator LEDs, the lamp indicator, the final lamp, and optional multimeter readings.

The system is intentionally simple, but it captures properties that are important in real diagnostic settings. Faults can propagate through the chain, so an upstream fault can cause downstream indicators and the final lamp to turn off. Different faults can produce similar observations, such as a detached switch, broken cable, or missing upstream voltage. Some observations may also be unavailable, for example when indicator LEDs are removed or no multimeter is available.



\subsection{Bayesian Network Reference Model}
\label{sec:bn_reference_model}

A Bayesian network reference model was developed collaboratively as part of the group project and is used here as a reference structure for the belief-network implementation. The model was built from the physical logic of the ten-module electrical system: fault nodes influence voltage propagation, voltage states influence indicator and lamp observations, and optional multimeter readings provide additional evidence. The full BN specification, including graph structure, priors, CPTs, diagnostic scenarios, and implementation files, is documented in the shared project repository \cite{BNRepository2026}.

The numerical probabilities in the BN should not be interpreted as real industrial reliability data. Most conditional relations are deterministic because they follow the circuit logic. For the fault priors and the few non-deterministic relations, a large language model was used as a simulated domain expert to propose plausible and internally consistent values. This is reported as part of the modelling procedure, not as empirical reliability evidence. The values are intended to create a realistic toy-domain reference model that produces meaningful diagnostic ambiguity, while the thesis focuses on what happens when such probabilistic information is incomplete, unavailable, or weakly justified. This limitation is consistent with the broader challenge of expert probability elicitation, where subjective judgements can introduce bias or overconfidence \cite{OHagan2006}.

\subsection{Diagnostic Scenarios and Scope}
\label{sec:diagnostic_scenarios_scope}

The diagnostic scenarios are used to evaluate model behaviour under different evidence patterns. Some scenarios are relatively clear, such as a lamp fault where upstream indicators function but the final lamp is off. Other scenarios are more ambiguous, such as an upstream interruption where several downstream indicators are off. Some scenarios remove observations, for example by assuming that control-module LEDs are unavailable or that no multimeter measurement can be taken. The scenarios do not define the model structure; they are used to test the behaviour of a model designed from the physical system.

Several simplifying assumptions are made to keep the project manageable. The model does not include all possible electrical faults. Diode faults, LED reliability faults, internal wiring faults, and detailed polarity faults are outside the selected scope. The system is also discretised: voltage is represented using states such as 12V and 0V, and observations are represented as on or off. This abstraction is appropriate because the objective is not electrical simulation, but the representation of diagnostic uncertainty under incomplete probabilistic information.


\section{Building the Belief Network}
\label{sec:building_belief_network}

This section summarises the implementation workflow. Its purpose is practical: to show how a known diagnostic structure can be represented as a belief network while preserving missing information instead of completing it artificially.

\subsection{Variables, Relations, and Evidence}
\label{sec:variables_relations_evidence}

The first step is to define discrete variables. Fault variables represent possible root causes, voltage variables represent internal propagation states, and observable variables represent LED states, lamp states, and measurements. The second step is to define relations between variables. Deterministic physical relations, such as voltage propagation through a connected cable and closed switch, are encoded as belief valuations that assign mass to compatible configurations. Conditional diagnostic relations are encoded as conditional valuations; when the relation is fully known, the valuation behaves like a Bayesian conditional table, while partially known relations can leave mass on broader sets of possibilities.

Evidence is represented by assigning support to observed states. Observing that a given LED is off, for example, creates evidence for the corresponding observable variable. Evidence does not directly identify the root fault; it is combined with the model's structural relations and propagated through the network to obtain belief and plausibility values for fault variables.

\subsection{Representing Missing Information}
\label{sec:representing_missing_information}

The central implementation problem is how to represent missing probabilities. In a BN, a missing prior or CPT row must be filled before standard inference can run. In a belief network, missing information can instead be represented as ignorance. For example, if the prior probability of a fault variable is unknown, the belief network can assign mass to the full frame of that variable. This states that the model has no support favouring one state over another. Similarly, if a conditional probability row is missing for a parent configuration, the belief-network representation can avoid inventing a precise distribution for that row.

This treatment is the main reason belief networks may be useful for industrial users. In early-stage systems, rare-fault settings, or expert-supported diagnosis, the structure of a system may be known before reliable probability estimates are available. A belief network allows the modeller to use the available structural knowledge while keeping incomplete quantitative knowledge visible.

\subsection{Implementation with the Belief Function Machine}
\label{sec:implementation_bfm}

The implementation uses BFM in MATLAB \cite{GiangShenoy2003BFM}. Variables, relations, conditional relations, valuations, and conditional valuations are written in a Unified Input Language file. The file is translated into BFM format, loaded in MATLAB, and conditional beliefs are converted into unconditional beliefs using Smets' conditional embedding. Inference is then performed by combining and marginalising belief functions. The detailed command-level workflow is given in Appendix~\ref{app:additional_belief_network_details}. The important conceptual step is the conditional embedding, which allows Bayesian-style conditional knowledge to be represented inside the belief-function framework \cite{JirousekKratochvilShenoy2023}.


\section{Experimental Design and Results}
\label{sec:experimental_results}

This section evaluates the belief-network implementation. First, the complete belief-function model is validated against the complete BN reference model. Second, information is progressively removed to test whether belief-network inference remains diagnostically useful when priors, measurements, or observation relations are missing or weakened. The goal is not to show that belief networks are more accurate than BNs under complete information. The goal is to test whether the belief-network output remains useful when precise probabilistic modelling is no longer justified.

\begin{table}[H]
\caption{Experimental levels used in the belief-network evaluation.}\label{tab:experiment_levels}
\centering
\small
\begin{tabularx}{\textwidth}{p{1.0cm}p{3.0cm}X}
\toprule
Level & Information condition & Purpose \\
\midrule
L0 & Complete BFM model & Validate that the belief-network implementation reproduces the complete BN under precise information. \\
L1 & Local prior ignorance & Replace only the locally relevant fault priors by ignorance. \\
L2 & Fault-family prior ignorance & Replace broader relevant prior families, such as switch/cable priors or power-subsystem priors, by ignorance. \\
L4 & Combined incompleteness & Combine missing measurements, missing prior knowledge, and weakened observation reliability. \\
\bottomrule
\end{tabularx}
\end{table}

The numbering follows the experiment labels used in the implementation. An intermediate observation-reliability condition was used as a robustness check, but the main text reports L4 because it combines that weakening with missing priors and missing measurements. The results are interpreted using belief, plausibility, width ($Pl-Bel$), diagnostic ranking, localisation, and empty-set mass. Width is important because it shows where support is unresolved rather than committed to one fault.

\subsection{Complete-information validation: L0}
\label{sec:l0_validation}

The L0 model was used as a validation bridge between the complete BN and the BFM implementation. In this condition, all priors and diagnostic relations from the BN were encoded in the BFM model without additional ignorance. Belief and plausibility collapsed to the same value for all queried singleton fault states, so the L0 belief-function model can be compared directly with the BN posteriors.

\begin{table}[H]
\caption{L0 validation results under complete information.}\label{tab:l0_results}
\centering
\small
\begin{tabularx}{\textwidth}{p{0.8cm}p{2.8cm}p{4.1cm}X}
\toprule
Scen. & Diagnostic case & Main L0 output & Interpretation \\
\midrule
7 & Battery exhaustion & \texttt{Batt = exhausted}: Bel.=Pl.=1.00 & BFM reproduces the complete BN conclusion. \\
8 & Module-3 interruption & \texttt{SW3 = detached}: 0.607; \texttt{C3 = broken}: 0.405 & BFM reproduces the BN ranking and localises the fault around module 3. \\
11 & Sparse downstream failure & Support distributed across several switch/cable faults & Sparse evidence already prevents precise localisation under complete information. \\
14 & PSU short with battery discharge & \texttt{PSF = yes}: 1.00; \texttt{Batt = exhausted}: 1.00 & BFM reproduces the compound upstream diagnosis. \\
\bottomrule
\end{tabularx}
\end{table}

These results show that, when all information is precise, the BFM implementation is a faithful translation of the BN for the tested scenarios. Scenario 14 is not contradictory: \texttt{PSF} and \texttt{Batt} are different fault variables and can both be true because a PSU short can discharge the battery.

\subsection{Local prior ignorance: L1}
\label{sec:l1_local_prior}

The L1 condition tested local prior ignorance. Only scenario-specific prior information was replaced by ignorance, while the diagnostic structure, observations, and deterministic relations remained available. This represents a realistic situation in which a company understands the system but no longer trusts a specific component failure-rate estimate, for example after redesign or supplier change.

\begin{table}[H]
\caption{L1 results under local prior ignorance.}\label{tab:l1_results}
\centering
\small
\begin{tabularx}{\textwidth}{p{0.8cm}p{2.9cm}p{3.8cm}p{1.1cm}X}
\toprule
Scen. & Missing information & Main L1 output & Width & Interpretation \\
\midrule
7 & Battery-related local prior & \texttt{Batt}: Bel.=Pl.=1.00 & 0.00 & Direct battery evidence preserves the diagnosis. \\
8 & Local priors for \texttt{SW3} and \texttt{C3} & \texttt{SW3}, \texttt{C3}: Bel.=0, Pl.=1.00 & 1.00 & Both local explanations remain possible, but the model refuses to rank them precisely. \\
11 & Local priors for \texttt{SW6} and \texttt{C6} & \texttt{SW6}, \texttt{C6}: Bel.=0, Pl.=1.00 & 1.00 & Relevant local faults remain possible, but sparse evidence prevents precise commitment. \\
14 & PSU-short and battery-related local priors & \texttt{PSF}, \texttt{Batt}: Bel.=Pl.=1.00 & 0.00 & Direct measurements dominate missing local prior information. \\
\bottomrule
\end{tabularx}
\end{table}

Scenario 8 is the clearest L1 result. Under complete information, the BN and L0 BFM prefer \texttt{SW3 = detached} over \texttt{C3 = broken}, mainly because of their relative priors. When those priors are removed, both faults receive belief 0 and plausibility 1.00. This is not a failure: the model still localises the fault to the module-3 transition, but it no longer claims a justified subcomponent-level ranking.

\subsection{Fault-family prior ignorance: L2}
\label{sec:l2_family_prior}

The L2 condition removed broader prior families. In Scenarios 8 and 11, all switch and cable priors were replaced by ignorance. In Scenarios 7 and 14, the upstream power-family priors were made ignorant. This is a stronger test because an entire reliability database is unavailable, not just one local prior.

\begin{table}[H]
\caption{L2 results under fault-family prior ignorance.}\label{tab:l2_results}
\centering
\small
\begin{tabularx}{\textwidth}{p{0.8cm}p{2.6cm}p{4.1cm}p{1.1cm}X}
\toprule
Scen. & Missing information & Main L2 output & Width & Interpretation \\
\midrule
7 & Upstream power-family priors & \texttt{Batt}: Bel.=Pl.=1.00 & 0.00 & Direct evidence still identifies battery exhaustion. \\
8 & All switch/cable priors & Downstream switch/cable faults remain plausible; upstream faults ruled out & 1.00 & Diagnosis broadens, but remains physically constrained to the downstream region from module 3. \\
11 & All switch/cable priors with sparse evidence & Broad plausibility over switch/cable candidates & up to 1.00 & The model becomes cautious because neither observations nor priors justify exact localisation. \\
14 & Upstream power-family priors & \texttt{PSF}, \texttt{Batt}: Bel.=Pl.=1.00 & 0.00 & Direct PSU-short and battery measurements preserve the compound diagnosis. \\
\bottomrule
\end{tabularx}
\end{table}

In Scenario 8, L2 widens the result compared with L1: the model no longer distinguishes sharply among the downstream switch/cable family. However, the uncertainty is still structured. Faults before module 3 are ruled out because the evidence shows that power reached modules 1 and 2. Thus, the model does not become randomly ignorant; it preserves the physical constraints imposed by the remaining evidence.

\subsection{Combined incompleteness: L4}
\label{sec:l4_combined}

The L4 condition tested more realistic mixed incompleteness. Scenario 15 removed the PSU-short measurement from a battery-discharge case and weakened the power-subsystem priors and visual rules. Scenario 16 removed the boundary indicator \texttt{I3} from the module-3 interruption case, while also making switch/cable priors ignorant and weakening visual observation rules.

\begin{table}[H]
\caption{L4 results under combined incompleteness.}\label{tab:l4_results}
\centering
\small
\begin{tabularx}{\textwidth}{p{0.8cm}p{3.0cm}p{4.3cm}p{1.1cm}X}
\toprule
Scen. & Combined incompleteness & Main L4 output & Width & Interpretation \\
\midrule
15 & PSU-short measurement unavailable; power priors ignorant; visual rules weakened & \texttt{Batt}: Bel.=Pl.=1.00; \texttt{PSF}: Bel.=0, Pl.=1.00 & 0.00 / 1.00 & Battery exhaustion is directly supported, while hidden PSU short remains unresolved but plausible. \\
16 & \texttt{I3} unavailable; switch/cable priors ignorant; visual rules weakened & \texttt{SW3--SW8}, \texttt{C3--C8}, and cable-to-load: Bel.=0, Pl.=1.00 & 1.00 & Exact localisation is lost, but the downstream fault region remains fully plausible. PSU short and battery are excluded; modules 1--2 retain only low residual plausibility. \\
\bottomrule
\end{tabularx}
\end{table}

Scenario 15 shows the value of separating belief from plausibility. Battery exhaustion is directly supported by the 0V battery measurement, while a hidden PSU short receives no belief but remains fully plausible because the measurement that would confirm or reject it is missing. Scenario 16 shows graded degradation. The power-supply short and battery exhaustion are excluded, while early switch/cable faults retain only low residual plausibility (approximately 0.04--0.20) because the visual rules were weakened. The model therefore does not rule out every upstream component in a binary way; instead, it strongly downweights them while keeping the downstream region as the main diagnostic envelope.

\subsection{Conflict and interpretation}
\label{sec:conflict_interpretation}

The empty-set mass, $m(\emptyset)$, was retained as a separate indicator of conflict between the combined evidence and the model assumptions. In these experiments, conflict mainly reflects prior-vs-evidence tension rather than sensor disagreement: precise priors encode faults as rare, while the observed evidence indicates that a fault has occurred. The BN normalises this tension away when conditioning on evidence, while the BFM exposes it as conflict.

\begin{table}[H]
\caption{Empty-set mass across experimental levels.}\label{tab:conflict_results}
\centering
\small
\begin{tabular}{ccccc}
\toprule
Scenario & L0 & L1 & L2 & L4 \\
\midrule
7  & 0.920 & 0.005 & 0.000 & -- \\
8  & 0.959 & 0.173 & 0.085 & -- \\
11 & 0.683 & 0.084 & 0.084 & -- \\
14 & 0.995 & 0.000 & 0.000 & -- \\
15 & --    & --    & --    & 0.000 \\
16 & --    & --    & --    & 0.085 \\
\bottomrule
\end{tabular}
\end{table}

The experiments show three patterns. First, decisive evidence can preserve the correct diagnosis even when priors are weakened, as in Scenarios 7 and 14. Second, when evidence localises a fault region but does not justify a precise component distinction, as in Scenario 8, the BFM preserves the relevant candidates as plausible rather than inventing a precise ranking. Third, when both observations and priors are weak, as in Scenarios 11 and 16, the model becomes broad but remains physically constrained. This supports the practical interpretation of belief networks: they are useful not because they are more precise than BNs, but because they show when precision is no longer justified.

\section{Practical Applicability and Limitations}
\label{sec:practical_guidance}

\subsection{When Belief Networks Are Recommended}
\label{sec:when_recommended}

Belief networks are recommended when the diagnostic structure is known well enough to model, but precise probabilities are incomplete, weakly justified, or unavailable. This can occur in early-stage product development, rare-fault diagnosis, safety-critical systems with limited failure data, or situations where expert knowledge is qualitative rather than numerical. The experiments suggest that the method is most useful when the goal is not only to rank faults, but also to see which alternatives have truly been ruled out.

A practical recommendation is therefore: use a belief network when preserving ignorance is part of the diagnostic task. If missing probabilities would otherwise be replaced by arbitrary priors, a belief network can communicate a more cautious result. It may still identify a fault directly when evidence is decisive, but when evidence is incomplete it produces a diagnostic envelope: supported faults, plausible but unresolved faults, and excluded faults.

\subsection{When Belief Networks Are Not Recommended}
\label{sec:when_not_recommended}

Belief networks are not always preferable. If sufficient data or reliable expert estimates are available for the priors and CPTs, a BN is usually simpler, faster, easier to communicate, and better supported by software. This is consistent with comparative work showing that Bayesian reasoning is often preferable when information is sufficiently complete and non-conflicting, while Dempster--Shafer reasoning becomes more informative when incompleteness or conflict is central \cite{VerbertBabuskaDeSchutter2017}.

The method is also risky when evidence sources are not distinct. Dempster's rule should be applied to distinct bodies of evidence; otherwise, repeated use of the same underlying information can cause overconfidence \cite{Shenoy2023DistinctBeliefFunctions}. This matters in diagnosis because several observations may share the same physical dependency or sensor chain. In such cases, cautious combination strategies or a redesigned evidence model may be needed \cite{Denoeux2008}.

\subsection{Scalability, Tooling, and Technology Readiness}
\label{sec:scalability_tooling_readiness}

The main technical limitation is computational complexity. Because belief functions can assign mass to sets of hypotheses, the number of focal elements can grow quickly, especially in larger networks with many variables and states. Practical deployment would require modular modelling, local computation, and possibly approximation methods \cite{ShaferShenoyMellouli1987,Shenoy2023IncompleteBN}.

Software maturity is another limitation. Bayesian networks have a larger ecosystem of tools, libraries, tutorials, and industrial examples. BFM provides a useful research environment, but it is not a production-ready industrial platform. MATLAB should therefore be understood as the research platform used in this thesis, not as the only possible implementation environment for the underlying method.

From a technology-readiness perspective, the method should be treated as research-prototype level rather than deployment-ready. TRLs provide a way to discuss maturity consistently \cite{Mankins1995}. The underlying theory is mature, and research implementations exist, but industrial deployment would still require validation on realistic systems, scalable implementation, careful treatment of evidence dependence, and user-facing explanations of belief and plausibility.

\subsection{Future Work}
\label{sec:future_work}

Future work should test the approach on larger diagnostic systems, compare alternative representations of partial conditional knowledge, and investigate approximation strategies for larger belief networks. More work is also needed on interfaces that explain not only which faults are plausible, but why the model remains uncertain and which additional measurements would reduce that uncertainty. Finally, future implementations could explore Python, Java, or C++ environments to separate the practical method from the MATLAB-based research prototype.


\section{Conclusion}
\label{sec:conclusion}

This thesis investigated when belief networks are suitable for fault diagnosis under incomplete probabilistic information. The modular electrical case study showed that a complete belief-function model can reproduce the Bayesian reference model under complete information, while progressively incomplete versions reveal what happens when priors, measurements, or observation relations are no longer fully trusted. The results support a practical conclusion: belief networks are useful when the diagnostic structure is known but precise probabilities are weakly justified, because they preserve unresolved alternatives through belief and plausibility instead of forcing precise posterior rankings. They are not a universal replacement for Bayesian networks. When reliable probabilities are available, BNs remain simpler and more mature. Belief networks are most appropriate as an expert-supported diagnostic approach when missing probabilistic information is itself an important part of the problem, and when users can benefit from seeing which faults are supported, which remain plausible, and which have genuinely been ruled out.


\appendix

\section{Additional Dataset Details}
\label{app:additional_dataset_details}

The main text reports only representative diagnostic scenarios. The full scenario catalogue and evidence files are maintained in the project repository and accompanying CSV outputs. The selected scenarios cover decisive upstream faults, mid-chain interruptions, sparse observations, compound upstream faults, and derived missing-measurement cases.

\section{Additional Bayesian Network Details}
\label{app:additional_bn_details}

The full Bayesian network specification, including the DAG, priors, CPTs, and implementation notebooks, is available in the shared repository \cite{BNRepository2026}. The model was used as a complete-information reference and not as a claim of industrial reliability calibration.

\section{Additional Belief Network Details}
\label{app:additional_belief_network_details}

The BFM workflow used in the experiments follows the standard sequence documented for the tool: write variables, relations, valuations, and conditional valuations in UIL; translate the file using \texttt{uil2bm}; load the generated model; convert conditional beliefs using \texttt{condiembed}; retain the embedded beliefs with \texttt{keepbel}; and query fault variables using \texttt{solve} and \texttt{showbel} \cite{GiangShenoy2003BFM}.

\section{Additional Results}
\label{app:additional_results}

Extended CSV outputs for L0, L1, L2, and L4 contain all queried fault candidates, belief, plausibility, width, empty-set mass, and rankings. The main paper reports only the values needed to support the experimental argument.


% ---- Bibliography ----
\bibliographystyle{splncs04}
\bibliography{references_review_implemented}

\end{document}
